% =====================================================================
% ICSE 2027 NIER -- PUBLIC version for arXiv (double-anonymity removed)
% Generated 2026-09-02 from nier-icse2027.tex
% ---------------------------------------------------------------------
% USAGE (arXiv, 2026 endorsement policy):
%   1. Register arXiv account (real name + ORCID).
%   2. Get endorsement for cs.SE:
%        - auto: institutional email + existing claimed papers, OR
%        - manual: contact an eligible endorser from your cited papers
%          (see docs/paper/submission-checklist.md appendix F for the
%          candidate list: Abrahamsson / Waseem / Tang / Fan / Kemell).
%   3. Post to arXiv AFTER the ICSE 2027 NIER notification (2026-12-18);
%      never post "under submission" during double-anonymous review.
%   4. Upload .tex + .bib only -- arXiv compiles them (do NOT upload PDF).
%   5. Subjects: cs.SE (primary), cs.AI (secondary cross-list).
%   6. License: arXiv default "perpetual, non-exclusive license" or
%      CC BY 4.0 -- publish BEFORE signing the ACM copyright form.
%   7. Do not write "under submission" in the arXiv Comments field;
%      after acceptance you may write "To appear at ICSE 2027 (NIER)".
% Compile: pdflatex nier-icse2027-public && bibtex nier-icse2027-public && pdflatex nier-icse2027-public && pdflatex nier-icse2027-public
% =====================================================================
\documentclass[sigconf]{acmart}

% --- real author block (FILL IN before posting) ---------------------
\author{First Author}
\affiliation{%
  \institution{Your Institution}
  \city{Your City}
  \country{Your Country}}
\email{first.author@example.com}
\orcid{0000-0000-0000-0000}

\author{Second Author}
\affiliation{%
  \institution{Your Institution}
  \city{Your City}
  \country{Your Country}}
\email{second.author@example.com}
\orcid{0000-0000-0000-0000}

% --- conference metadata ---------------------------------------------
\acmConference[ICSE '27]{International Conference on Software Engineering}{2027}{TBD}
\acmYear{2027}
\copyrightyear{2027}

\begin{document}

\title{Agentic Metabolic Engineering: Managing Workspace Byproducts in AI-Driven Software Development}

\begin{abstract}
The rapid adoption of vibe coding and agentic software development has created a new class of technical debt: the uncontrolled accumulation of workspace byproducts---caches, logs, draft files, and abandoned implementations. While recent work has identified workspace state drift as a fundamental bottleneck in agent performance, no existing framework provides a systematic, policy-driven approach to managing the lifecycle of these byproducts. We introduce \textbf{Agentic Metabolic Engineering}, a paradigm that treats the agent workspace as a living system requiring digestion, not just deletion. We define four metabolic phases---catabolism (audit), sequestration (clean), verification (verify), and anabolism (rollback)---and present \texttt{workspace-metabolism}, a zero-dependency Python CLI that implements this paradigm. Through a reproducible benchmark simulating 30 agentic loops, we show that a metabolized workspace maintains 2 active files versus 242 in an unmanaged workspace, while preserving full auditability and rollback capability. Our framework fills the gap between workspace governance during execution (Harness Engineering, Loop Engineering) and workspace cleanup after execution, offering the first verifiable, policy-driven lifecycle management for agent-generated byproducts.
\end{abstract}

\begin{CCSXML}
<ccs2012>
<concept>
<concept_id>10011007.10011006.10011008</concept_id>
<concept_desc>Software and its engineering~Software creation and management</concept_desc>
<concept_significance>500</concept_significance>
</concept>
<concept>
<concept_id>10011007.10011006.10011041</concept_id>
<concept_desc>Software and its engineering~Software maintenance tools</concept_desc>
<concept_significance>500</concept_significance>
</concept>
</ccs2012>
\end{CCSXML}
\ccsdesc[500]{Software and its engineering~Software creation and management}
\ccsdesc[500]{Software and its engineering~Software maintenance tools}

\keywords{agentic software engineering, technical debt, workspace management, vibe coding, file lifecycle}

\maketitle

\section{Introduction}
The emergence of vibe coding---software development through natural language prompts to AI agents---has fundamentally changed what a software workspace looks like. As of 2026, 92\% of US-based developers (76\% globally) use AI coding tools daily \cite{jetbrains2026}, and 41\% of all code produced globally is AI-generated. While 93.5\% of technology leaders rate AI-generated code as higher quality at review time, 78\% report more production incidents once this code ships \cite{newrelic2026}.

But code quality is only part of the problem. Each agentic loop---plan, execute, observe, reflect---leaves behind files: caches, logs, draft outputs, half-finished refactors, duplicated lock files. The workspace becomes a compost pile without a gardener.

Recent work has begun to recognize this problem. LemonHarness identified that ``agents typically observe only tool outputs and log fragments, while the actual state changes occur in the file system'' \cite{ren2026lemonharness}. SABER shows that failures in long-horizon agents ``cluster at a small slice of mutating steps,'' such as file deletion \cite{saber2025}. Workspace-Bench~1.0 established a benchmark for evaluating agents on workspace tasks with large-scale file dependencies \cite{tang2026workspacebench}. AgentFold treats agent context as ``a dynamic cognitive workspace to be actively sculpted'' \cite{tongyi2025agentfold}.

These works establish that \textbf{workspace management is not a minor hygiene issue but a fundamental bottleneck} in agent performance. Yet they stop at governance \emph{during} execution or measurement \emph{after} execution. None provides a systematic framework for managing the byproducts that accumulate \emph{between} loops.

\textbf{Our contribution} is \textbf{Agentic Metabolic Engineering}---a paradigm that treats the agent workspace as a living system with four lifecycle phases, implemented in \texttt{workspace-metabolism}, a zero-dependency Python CLI.

\section{The Problem: Workspace Byproducts}
Vibe coding creates a characteristic loop: AI generates code $\rightarrow$ errors occur $\rightarrow$ AI fixes them $\rightarrow$ new dependencies appear $\rightarrow$ previous approaches are abandoned but not removed $\rightarrow$ the cycle repeats. The result: \texttt{archive/}, \texttt{deprecated/}, \texttt{test\_drafts/}, \texttt{tmp/}, half-finished refactors grow \textbf{faster than the code itself}.

Existing approaches are insufficient:
\begin{itemize}
\item \textbf{Direct deletion (\texttt{rm -rf})} is \textbf{autophagy}---it burns potentially valuable intermediate matter.
\item \textbf{Git} is right for tracked source, but most byproducts are untracked, high-churn, non-code assets.
\item \textbf{Runtime hygiene tools} like \texttt{zclean} \cite{zclean2026} and \texttt{agent-gc} \cite{jeong2026agentgc} target processes or caches, but lack policy-driven classification, auditability, and rollback.
\end{itemize}

What is missing is a framework that treats workspace byproducts as having a \textbf{lifecycle}---not as waste to be blindly deleted, but as matter to be digested, tracked, and potentially reabsorbed.

\section{Agentic Metabolic Engineering}
\textbf{Agentic Metabolic Engineering} is the systematic management of byproducts generated by agent-driven software workspaces. Where \textbf{Harness Engineering} \cite{harness2026} asks ``Can the agent complete this task reliably?'' and \textbf{Loop Engineering} \cite{loops2026} asks ``Can the agent run itself without us?'', \textbf{Agentic Metabolic Engineering} asks: ``After the agent completes tasks, can the workspace survive the next one?'' We position it as the fifth layer of the agentic engineering stack---after Prompt Engineering \cite{schulhoff2024}, Context Engineering \cite{context2025}, Harness Engineering \cite{harness2026}, and Loop Engineering \cite{loops2026}---the last mile of every loop: after the agent plans, executes, observes, and reflects, something must tend to the files it leaves behind.

\subsection{The Four Phases}
Each phase maps to a command in \texttt{workspace-metabolism}:
\begin{table}[h]
\caption{The four metabolic phases and their commands.}
\label{tab:phases}
\begin{tabular}{lll}
\textbf{Phase} & \textbf{Command} & \textbf{Description}\\
\hline
Catabolism & \texttt{wm audit} & Scan workspace, identify candidates\\
Sequestration & \texttt{wm clean} & Move expired matter to recycle area\\
Verification & \texttt{wm verify} & Hash-chain journal detects tampering\\
Anabolism & \texttt{wm rollback} & Re-inject material back into workspace\\
\end{tabular}
\end{table}

\subsection{A Formal Framework}
We define a three-dimensional framework:
\begin{itemize}
\item \textbf{Metabolism Ladder (M0--M3)}: M0 = audit only; M1 = audit + dry-run plan; M2 = audit + clean (move to recycle, SHA-256 per file); M3 = full closed loop (audit + clean + verify + rollback). Most cleanup tools operate at M1 or below; our tool targets M3.
\item \textbf{Policy Hierarchy (G1--G4)}: G1 = never touch; G2 = monitor; G3 = require approval; G4 = auto-process.
\item \textbf{Metabolic Cycle}: Trigger (scheduled or event-driven) $\rightarrow$ Audit $\rightarrow$ Clean $\rightarrow$ Verify $\rightarrow$ (optionally) Rollback $\rightarrow$ Outcome state (active workspace + recycle area + audit journal).
\end{itemize}

\subsection{Policy as Code}
A single JSON policy file grades every path, enabling versionable, reviewable, auditable governance.

\subsection{The Micro-Metabolism Pattern}
At the end of each agentic loop, a micro-metabolism step executes: ``The files just produced---keep, archive, recycle, or leave for tomorrow's digestion?'' This can be a wrapper around \texttt{wm audit --json} that an agent calls before its next planning phase.

\section{Implementation \& Evaluation}
\texttt{workspace-metabolism} is a zero-dependency Python CLI (3.11+) implementing the paradigm. All operations default to dry-run; destructive actions require explicit flags.

\textbf{Safety Model:} (1) Never delete during clean---files are moved; (2) Cryptographic verification---SHA-256 per file; (3) Hash-chain audit---append-only, tamper-evident; (4) Explicit purge---real deletion only in recycle area.

\textbf{Benchmark.} We constructed a reproducible benchmark simulating 30 agentic loops on two identical workspaces: one with \texttt{wm clean} after each loop, one without. Each loop generates 8 byproducts typical of AI coding sessions. The benchmark uses a fixed random seed and is checked into the repository as \texttt{examples/metabolism\_benchmark.py}.

\textbf{Results.} After 30 loops:
\begin{table}[h]
\caption{Benchmark results after 30 simulated agentic loops.}
\label{tab:results}
\begin{tabular}{lcc}
\textbf{Metric} & \textbf{With Metabolism} & \textbf{Without Metabolism}\\
\hline
Active files & \textbf{2} & 242\\
Recycle area & 240 files, verified & 0\\
Hash-chain integrity & $\checkmark$ Passed & N/A\\
Rollback capability & $\checkmark$ Verified & N/A\\
\end{tabular}
\end{table}

\textbf{2 vs.\ 242.} The metabolized workspace kept exactly what it needed, moved everything else to a recyclable area with cryptographic verification, and maintained a tamper-evident audit trail; recycled files are recoverable via \texttt{wm rollback <run\_id>} (verified at the byte level). The unmanaged workspace accumulated 240 byproducts.

\section{Related Work}
\textbf{Vibe Coding \& Technical Debt.} Waseem et al.\ analyzed ``flow-debt tradeoffs'' in vibe coding \cite{waseem2025}. A 2025 survey identified dependency proliferation as ``a structural form of technical debt'' \cite{vibesurvey2025}. The position paper ``Vibe Coding Needs Vibe Reasoning'' calls out technical debt accumulation as a key limitation \cite{position2025}.

\textbf{Agent Workspace Governance.} LemonHarness establishes explicit execution boundaries \cite{ren2026lemonharness}. AgentFold treats context as a dynamic cognitive workspace \cite{tongyi2025agentfold}. Workspace-Bench~1.0 benchmarks agents on workspace tasks \cite{tang2026workspacebench}. SABER shows failures concentrate at mutating actions \cite{saber2025}.

\textbf{Workspace Hygiene Tools.} \texttt{zclean} \cite{zclean2026} and \texttt{agent-gc} \cite{jeong2026agentgc} target orphaned processes and build artifacts but lack policy-driven classification, auditability, and rollback.

\textbf{Our Position.} To our knowledge, \textbf{no existing work combines policy-driven classification, recyclable cleanup with cryptographic rollback, hash-chain auditability, and native agent integration} into a unified framework. \texttt{workspace-metabolism} is the first to treat workspace byproducts as a \textbf{lifecycle}---not as waste to be deleted, but as matter to be digested, tracked, and potentially reabsorbed.

\section{Future Plans}
This section outlines the work we plan to turn this new idea into a full-length paper.

\textbf{1. Real-World Validation.} Our benchmark uses simulated agentic loops. We plan to collect longitudinal data from real AI development workflows, measuring metabolic debt accumulation and its correlation with agent performance metrics (task completion time, error rates).

\textbf{2. Workspace Health Score.} Our shipped 0--100 health score (journal auditability, governance coverage, rot burden, recycle readiness) already enables CI/CD gating; we plan to decompose it into a multidimensional model aligned with data-quality frameworks, and to study cross-repository comparability.

\textbf{3. Deeper Agent Integration.} The shipped MCP server already lets agents invoke \texttt{wm audit --json} in their observation phase, and a governance proxy (\texttt{wm gate}) checks every tool call against policy; we plan to close the loop into fully autonomous micro-metabolism while preserving auditability.

\textbf{4. Metabolic Debt Theory.} We will formalize the relationship between metabolic debt and agent performance, developing predictive models of when and how unmanaged byproducts degrade agent effectiveness.

\section{Conclusion}
The rapid adoption of vibe coding has created a new challenge: the uncontrolled accumulation of workspace byproducts that degrade agent performance. We have introduced \textbf{Agentic Metabolic Engineering} as a paradigm for addressing this challenge, defined four metabolic phases, and presented \texttt{workspace-metabolism}, a zero-dependency Python CLI. Through a reproducible benchmark, we have shown that a metabolized workspace maintains order (2 active files) while an unmanaged workspace accumulates 242 files. If Harness Engineering is the skeleton of an agent and Loop Engineering its heartbeat, then \textbf{Agentic Metabolic Engineering is its digestive system}---keeping the workspace alive for the next cycle.

% Public artifact link (arXiv version only -- restore for camera-ready too):
The tool is open-source: \url{https://github.com/metabolism-tools/workspace-metabolism}

\begin{acks}
The authors thank the open-source community for feedback on early releases of \texttt{workspace-metabolism}.

% ACM/IEEE policy: disclose generative-AI assistance here (fill in truthfully).
% Example: "ChatGPT was utilized to generate sections of this Work,
% including text, tables, graphs, code, data, citations, etc."
\end{acks}

\bibliographystyle{ACM-Reference-Format}
\bibliography{nier-references}

\end{document}
